%! AP Statistics — Unit 1, Lesson 1.1 — Group Activity KEY
\documentclass[10pt]{article}
\usepackage{apstats-article}
\usepackage{apstats-key}

%=============================================================================
\begin{document}

\pageheader{Unit 1, Lesson 1.1}{Group Activity — Answer Key}
\begin{tabularx}{\linewidth}{X X X}
Name:\;\ans{Key} & Partner:\; & Period:\;
\end{tabularx}

\vspace{0.1in}

\begin{tcolorbox}[colback=sky, colframe=navy, boxrule=0.8pt, arc=2mm,
    left=3mm, right=3mm, top=2mm, bottom=2mm]
\small
For each headline below, identify the \textbf{individual}, \textbf{population},
\textbf{sample}, and \textbf{variable(s)}. Then answer the follow-up question.
Remember: the \emph{population} is who or what we \emph{want} to know about;
the \emph{sample} is who or what we \emph{actually studied}.
\end{tcolorbox}

\vspace{0.14in}

% ════════════════════════════════════════════════════════════════════════════
% SCENARIO 1 KEY
% ════════════════════════════════════════════════════════════════════════════
\begin{scenariobox}[Scenario 1]{navylight}

\begin{headlinebox}{sky}
\small\textit{``A survey of 500 U.S.\ high school students found that 78\% report
feeling stressed about grades at least once a week.''}
\end{headlinebox}

\vspace{8pt}
\componenttablekey
  {A U.S.\ high school student}
  {All U.S.\ high school students}
  {The 500 students surveyed}
  {How often the student feels stressed about grades (categorical / ordinal)}

\vspace{8pt}
\small\textbf{Would it be possible to conduct a census for this study?}\\[4pt]
\ans{No --- there are roughly 15 million high school students in the U.S.
Surveying all of them would be impossible due to cost, time, and logistics.
A well-chosen sample of 500 can still give reliable results.}

\end{scenariobox}

\vspace{0.12in}

% ════════════════════════════════════════════════════════════════════════════
% SCENARIO 2 KEY
% ════════════════════════════════════════════════════════════════════════════
\begin{scenariobox}[Scenario 2]{greenacc}

\begin{headlinebox}{greenbg}
\small\textit{``A national study tracked 1{,}200 elementary school students over
three years and found that those who had at least 30 minutes of recess daily
scored higher on standardized reading assessments.''}
\end{headlinebox}

\vspace{8pt}
\componenttablekey
  {An elementary school student}
  {All U.S.\ elementary school students (accept: all elementary school students)}
  {The 1{,}200 students tracked in the study}
  {Daily recess time (quantitative); score on standardized reading assessment (quantitative)}

\vspace{6pt}
\begin{teachernote}
\scriptsize
Students may define the population too narrowly (``students in this study'') or too broadly
(``all children in the world''). Push them to use the headline's scope --- ``national study''
implies U.S.\ students. Also worth noting: two variables are present; students who identify
only one are partially correct --- ask them what else was measured.
\end{teachernote}

\vspace{8pt}
\small\textbf{Does association mean causation?}\\[4pt]
\ans{No. The study shows that more recess \emph{goes along with} higher reading scores,
but there could be other explanations --- for example, schools that offer more recess
might also have better-funded programs overall, or students who get more exercise may
simply be more alert. To establish causation we would need a controlled experiment.}

\end{scenariobox}

\vspace{0.12in}

% ════════════════════════════════════════════════════════════════════════════
% SCENARIO 3 KEY
% ════════════════════════════════════════════════════════════════════════════
\begin{scenariobox}[Scenario 3 \normalfont\small\color{white!80!black}\quad stretch]{redacc}

\begin{headlinebox}{redbg}
\small\textit{``Scientists examined 60 coral reef sites across the Pacific Ocean
and found that water temperature had increased at 85\% of them over the past decade.''}
\end{headlinebox}

\vspace{8pt}
\componenttablekey
  {A coral reef site (not a person --- a location)}
  {All coral reef sites in the Pacific Ocean}
  {The 60 reef sites examined}
  {Water temperature change over the past decade (quantitative)}

\vspace{6pt}
\begin{teachernote}
\scriptsize
The most common error here is writing ``a scientist'' or ``the ocean'' as the individual.
Emphasize: the individual is whatever the researchers recorded data \emph{about} --- here,
each measurement belongs to a reef site. This is a high-value AP misconception to address early.
Population boundary is also genuinely ambiguous (Pacific only? all oceans?); accept either
as long as students justify their choice.
\end{teachernote}

\vspace{8pt}
\small\textbf{``The individual is the ocean.'' Do you agree?}\\[4pt]
\ans{No. The individual is a \emph{coral reef site} --- a specific, measurable location.
The ocean is the broader context, not the unit being observed. Each data point in the study
belongs to one reef site: its name (or coordinates) and its recorded temperature change.
The ocean as a whole was never measured as a single entity.}

\end{scenariobox}

\vspace{0.12in}

% ── Reflection ───────────────────────────────────────────────────────────────
\begin{tcolorbox}[colback=goldbg, colframe=goldacc, boxrule=0.8pt, arc=2mm,
    left=3mm, right=3mm, top=2mm, bottom=2mm]
\small\textbf{Reflection:} \quad Which scenario was hardest? What made it harder?\\[4pt]
\ans{Most students will say Scenario 3. The key difficulty is that the individual is not a
person --- it requires stepping back and asking ``what was each measurement \emph{about}?''
rather than ``who was in the study?'' This is exactly the habit AP free-response rewards.}
\end{tcolorbox}

\end{document}
